% =====================================================================
%  2bat-ud5-7.tex  —  SESSIÓ 7: La binomial amb el full de càlcul
%  (sense preàmbul: s'inclou des de main.tex)
% =====================================================================
\horatitol{Sessió 7 — La binomial amb el full de càlcul}%
{Treball tecnològic: calcular probabilitats binomials amb el full de càlcul, reproduir els casos fets a mà i representar la distribució en un gràfic de barres.}

\subsection{Idea clau}

Portem dues sessions calculant probabilitats binomials \textbf{a mà}. Avui deixem que el
\textbf{full de càlcul} faci el càlcul repetitiu, per centrar-nos a interpretar. Farem
dues coses: \textbf{reproduir} els casos ja resolts (per guanyar confiança) i, sobretot,
\textbf{representar} la distribució sencera en un gràfic de barres, que prepara el pas cap
a la corba normal.

\begin{idea}
\textbf{La funció binomial al full de càlcul.} La majoria de fulls de càlcul tenen una
funció del tipus \texttt{DISTR.BINOM.N} (o \texttt{BINOM.DIST}):
\[
\texttt{=DISTR.BINOM.N(k;\ n;\ p;\ acumulat)}
\]
\begin{itemize}[leftmargin=1.4em, itemsep=2pt]
  \item \texttt{k}: nombre d'èxits; \texttt{n}: repeticions; \texttt{p}: probabilitat
        d'èxit.
  \item \texttt{acumulat = FALS}: dona $P(\text{exactament }k)$.
  \item \texttt{acumulat = CERT}: dona $P(\text{com a màxim }k)$ (probabilitat acumulada).
\end{itemize}
\textbf{La regla d'or:} la màquina estalvia el càlcul, però cal saber \textbf{interpretar}
el que surt i comprovar que coincideix amb el fet a mà.
\end{idea}

\subsection{Exemples}

\begin{exemple}[reproduir un cas fet a mà]
A la sessió 5 vam calcular, per a $n=5$ i $p=0,6$, que $P(3)\approx 0,3456$. Al full de
càlcul escrivim
\[
\texttt{=DISTR.BINOM.N(3;\ 5;\ 0,6;\ FALS)}
\]
i obtenim $0,3456$: \textbf{coincideix}. Guanyem confiança que l'eina i el càlcul a mà
diuen el mateix.
\end{exemple}

\begin{exemple}[representar la distribució]
Per veure la distribució sencera de $B(5;0,6)$, posem els valors de $k$ ($0,1,2,3,4,5$)
en una columna i, al costat, \texttt{=DISTR.BINOM.N(k;\ 5;\ 0,6;\ FALS)}:
\begin{center}
\renewcommand{\arraystretch}{1.2}
\begin{tabular}{c c c c c c c}
\toprule
$k$ & 0 & 1 & 2 & 3 & 4 & 5 \\
\midrule
$P(k)$ & 0,010 & 0,077 & 0,230 & 0,346 & 0,259 & 0,078 \\
\bottomrule
\end{tabular}
\end{center}
En fer-ne un gràfic de barres es veu que el màxim és a $k=3$ (el valor esperat) i que les
barres cauen a banda i banda: una forma que, amb més repeticions, s'assemblarà a una
\textbf{campana}.
\end{exemple}

\subsection{Exercicis resolts}

\begin{exresolt}[triar bé l'argument «acumulat»]
Vols la probabilitat que en un grup de $8$ ($p=0,7$) en superin \textbf{exactament 6}.
Quina fórmula escrius i per què?
\solucio
Com que és \textbf{exactament} 6, l'argument acumulat ha de ser \texttt{FALS}:
\[
\texttt{=DISTR.BINOM.N(6;\ 8;\ 0,7;\ FALS)}\ \to\ 0,2965\approx 29,6\%.
\]
Si hi posessis \texttt{CERT}, obtindries $P(\text{com a màxim }6)$, una altra pregunta.
\resposta{\texttt{=DISTR.BINOM.N(6;8;0,7;FALS)} $\approx 29,6\%$.}
\end{exresolt}

\begin{exresolt}[usar l'acumulada]
Amb $n=10$ i $p=0,5$, com calcularies $P(\text{com a màxim }4)$ amb una sola fórmula?
\solucio
L'acumulada amb \texttt{CERT} dona directament $P(\text{com a màxim }k)$:
\[
\texttt{=DISTR.BINOM.N(4;\ 10;\ 0,5;\ CERT)}\ \to\ 0,377\approx 37,7\%.
\]
\resposta{\texttt{=DISTR.BINOM.N(4;10;0,5;CERT)} $\approx 37,7\%$.}
\end{exresolt}

\subsection{Exercicis per fer}

\noindent\textit{Full de pràctica tecnològica. Anota la fórmula i el resultat de cada
cas.}

\begin{exercicis}
\begin{exlist}
  \item Escriu la fórmula per calcular $P(\text{exactament }5)$ amb $n=10$ i $p=0,5$, i
        anota el resultat.
  \item Reprodueix al full de càlcul el cas de la sessió 6: $n=12$, $p=0,65$,
        $P(\text{exactament }8)$. Coincideix amb el $23,7\%$ calculat a mà?
  \item Amb $n=8$ i $p=0,7$, calcula $P(\text{com a màxim }5)$ amb \texttt{acumulat=CERT}.
  \item Construeix la taula de $B(6;0,5)$ per a $k=0,\dots,6$ i digues en quin $k$ és més
        alta la probabilitat.
  \item \textbf{(Com a mínim)} Com calcularies $P(\text{com a mínim }7)$ amb $n=10$,
        $p=0,7$, fent servir l'acumulada? (Pista: $P(k\ge 7)=1-P(\text{com a màxim }6)$.)
  \item \textbf{(Interpretació)} En representar $B(20;0,5)$ en un gràfic de barres, quina
        forma esperes veure? On estarà el punt més alt? Com es relaciona amb el valor
        esperat $\mu=n\cdot p$?
\end{exlist}
\end{exercicis}

% ---------------------------------------------------------------------
\fullsolucions{Sessió 7}
\begin{solucions}
\begin{sollist}
  \item \texttt{=DISTR.BINOM.N(5;10;0,5;FALS)} $\to 0,2461\approx 24,6\%$.
  \item \texttt{=DISTR.BINOM.N(8;12;0,65;FALS)} $\to 0,2367\approx 23,7\%$. Coincideix
        amb el càlcul a mà.
  \item \texttt{=DISTR.BINOM.N(5;8;0,7;CERT)} $\to 0,4482\approx 44,8\%$.
  \item Taula de $B(6;0,5)$: $P(0)=0,016$, $P(1)=0,094$, $P(2)=0,234$, $P(3)=0,313$,
        $P(4)=0,234$, $P(5)=0,094$, $P(6)=0,016$. El màxim és a $k=3$ (el valor esperat
        $\mu=6\cdot 0,5=3$).
  \item $P(k\ge 7)=1-\texttt{DISTR.BINOM.N(6;10;0,7;CERT)}=1-0,3504\approx 0,6496=65,0\%$.
  \item Una forma de \textbf{campana} simètrica, amb el punt més alt a $k=10$, que és
        justament el valor esperat $\mu=20\cdot 0,5=10$.
\end{sollist}
\end{solucions}
